\documentclass[10pt,a4paper]{article}
\usepackage{amsmath,amssymb,amsthm,graphicx}
\usepackage{xfrac}
\newenvironment{enumeratei}
{\bgroup\def\theenumi{\roman{enumi}}\def\theenumii{\arabic{enumii}}\begin{enumerate}}
{\end{enumerate}\egroup}

\hyphenation{addi-tion-al belong belongs denote exchan-ge exchan-ged exchan-ges}

\newcommand{\psfrac}[2]{\sfrac{(#1)}{#2}}

\newcommand\mto{\mathchoice{\longmapsto}{\mapsto}{\mapsto}{\mapsto}}
\let\oldbigsqcup\bigsqcup
\renewcommand{\bigsqcup}{\mathop{\textstyle\oldbigsqcup}\displaylimits}

\let\epsilon\varepsilon

\DeclareMathOperator{\Ess}{Ess}
\DeclareMathOperator{\Leb}{Leb}

\usepackage{tikz}
\usetikzlibrary{plotmarks}

\theoremstyle{plain}
\newtheorem{theorem}{Theorem}[section]
\newtheorem{lemma}[theorem]{Lemma}
\newtheorem{proposition}[theorem]{Proposition}
\newtheorem{corollary}[theorem]{Corollary}
\newtheorem{claim}{Claim}

\theoremstyle{definition}
\newtheorem{remark}[theorem]{Remark}
\newtheorem{example}[theorem]{Example}

\newcommand{\A}{\mathcal{A}}
\newcommand{\B}{\mathcal{B}}
\newcommand{\I}{\mathcal{I}}

\newcommand{\Z}{\mathbb{Z}}
\newcommand{\N}{\mathbb{N}}
\newcommand{\R}{\mathbb{R}}
\title{An ergodic interval exchange transformation has an infinite connection time}
\author{Przemyslaw Berk, Frank Trujillo}
\date{doi:10.5802/jep.302; CC BY4.0}
\begin{document}
\maketitle
\noindent\textit{Editorial note.} The counted statement is the author's labelled corollary environment at source0.tex lines 220-222; its proof environment opens at line 224, the line after the blank separator, so the pairing is by adjacency and the retained proof range 224-226 is the whole environment.\par\medskip
\noindent\textit{Editorial note.} The lemma lem: notallconnections (statement, source lines 191-193) is retained as context because the counted proof cites it at source line 225; it is the only outside label of the retained proof.\par\medskip
\noindent\textit{Editorial note.} This source is new to the material-authoring work and its preamble is built from the author's own macro and theorem block (source lines 8-45) rather than from a prepared adapter.\par\medskip
\noindent\textit{Editorial note.} This article carries no label-number dump and the retained counter values are left empty for the whole source, so the printed numbering of the card is the wrapper's default; the identity of the counted item is its author label cor: existenceofinfiniteorbit.\par\medskip
\noindent\textit{Editorial scope.} One original result of the pinned publisher article, counted once: existence of an infinite connection time for an ergodic interval exchange transformation (author label cor: existenceofinfiniteorbit).\par\medskip
\noindent\textit{Allowed standard prerequisites (editorial).}\begin{itemize}
\item Ergodicity of an interval exchange transformation. The transformation is ergodic with respect to Lebesgue measure, so it has no non-trivial invariant union of intervals.. It is the hypothesis of the statement.
\item Connection time. M(beta) is the first return time of the left endpoint of the subinterval beta to the set of discontinuities, possibly infinite.. It is the quantity whose finiteness is excluded.
\item Periodic behaviour from finite connection times. If all connection times except one are finite, the transformation has a finite decomposition into periodic components.. It is the contradiction derived in the proof.
\end{itemize}\medskip
\begin{lemma}\label{lem: notallconnections}
Let $T:I\to I$ be an interval exchange transformation given by a permutation $\pi$ and length vector $\lambda$. If for every $\beta\in \A\setminus \{\pi_0^{-1}(1)\}$ we have $M(\beta)<\infty$, then $T$ has only periodic orbits. {More precisely, the base interval $I$ can be decomposed in a finite number of periodic components, given by semi-closed intervals, such that the period is uniform on each of these components.}
\end{lemma}

\begin{corollary}\label{cor: existenceofinfiniteorbit}
Assume that $T$ is an ergodic IET. Then there exists {$\beta \in\mathcal A \setminus \{\pi_0^{-1}(1)\}$ such that $M(\beta) = +\infty$. }
\end{corollary}

\begin{proof}
Assume, for the sake of contradiction, that $T$ is ergodic but that the conclusion does not hold. Then, by~Lemma \ref{lem: notallconnections}, there exists a non-empty semi-closed interval $[a, b) \subseteq I$ and $N \geq 1$ such that $T^N\mid_{[a, b)}$ is the identity map in $[a, b)$. Therefore, the set $\bigcup_{i = 0}^{N - 1} T^i\big( \big[a, \psfrac{a + b}{2}\big]\big)$ is a non-trivial $T$-invariant set, which contradicts the ergodicity of $T$.
\end{proof}

\end{document}
